\section{Measurement implications and discussion}
\paragraph{Observe required behavior rather than implementation bookkeeping.}
The X05 and X28 corrections arose because alternative implementations could return a contract relevant result without populating an optional internal ledger. Applying the contract based rule to all 144 saved submissions in the three reviewed families changed four focal witness labels. This gives a practical rule for agent benchmarks: evaluate required outputs and externally visible effects, and use internal traces only as supporting evidence. X06 shows the complementary problem. Its executed witness establishes confidentiality for one mixed record, but does not establish delivery of useful public fields. Prospective X06 and F08 output rules are recorded in Appendix~\ref{app:observers} and do not alter the historical results.

\begin{table}[t]
\centering\footnotesize
\begin{tabularx}{\linewidth}{@{}p{2.9cm}XX@{}}
\toprule
Claim & Evidence & Alternative explanation retained\\
\midrule
Correct memory varies by configuration & C$-$N is positive, zero, and negative & Local harms can coexist with aggregate cancellation.\\
Reminder response varies by configuration & Codex contrasts and bounds agree in direction; MiniSWE does not & Interfaces and instructions may shape the response.\\
Visible statement differs from protection & Fifteen runs have protection without agreed prior assumption recognition & The codes do not identify memory use or a causal pathway.\\
Focal passes depend on the observer & X05/X28 corrections and separate X06 output requirements & Confidentiality can hold because delivery was rejected.\\
\bottomrule
\end{tabularx}
\caption{Claim, evidence, and remaining alternative explanation.}
\label{tab:claims}
\end{table}

\paragraph{Limitations.}
The controlled synthetic tasks provide explicit source assumptions and independently testable functionality and focal properties, but thirteen families with two repetitions do not establish external validity for repository scale software engineering. The families cover a constrained set of trust shifts rather than the full range of changes that can alter a procedure's applicability. The 52-run behavioral analysis is exploratory and based on two model assisted coding passes, with eight disagreements and no human validation. The experiment isolates supplied procedural memory under controlled target conditions; retrieval quality, naturally occurring repository histories, longer term memory dynamics, and model specific mechanisms require external validation.

\section{Conclusion}
Correct memory produced positive, zero, and negative changes in functional focal witness failures across the six configurations. Lower endpoint values sometimes reflected lost functionality, while matched traces showed both unchecked reuse and successful target adaptation. A generic boundary reminder reduced the endpoint in all three Codex configurations, including under the missing outcome bounds, but the MiniSWE configurations did not share that pattern. Correct procedural memory therefore cannot be treated as uniformly beneficial or harmful when applicability assumptions change. Evaluating transfer requires observing both task completion and the property that the original procedure implicitly relied on.
